\Legend
直接按题意维护已访问格点，能模拟少量飞行；但 $n$ 可达 $10^9$。

关键是每四次指令包含东、西、南、北各一次。把 $24$ 种 $P$ 按平面旋转和反射归成三种代表：\tt{NEWS}、\tt{NESW}、\tt{NSEW}。设 $q\ge0$，下表依次给出第 $4q+1$ 到 $4q+4$ 次飞行后的坐标：

| 指令序列 | $v_{4q+1}$ | $v_{4q+2}$ | $v_{4q+3}$ | $v_{4q+4}$ |
| --- | --- | --- | --- | --- |
| \tt{NEWS} | $(-q,q+1)$ | $(1-q,q+1)$ | $(-q-1,q+1)$ | $(-q-1,q)$ |
| \tt{NESW} | $(-q,1)$ | $(q+1,1)$ | $(q+1,0)$ | $(-q-1,0)$ |
| \tt{NSEW} | $(-q,1-q)$ | $(-q,-q-1)$ | $(1-q,-q-1)$ | $(-q-1,-q-1)$ |

这些式子在 $q=0$ 时可直接模拟得到。假定前 $q$ 块成立，逐个方向查看该行或列已有格点的区间端点，即可确认下一块四个坐标分别是该方向第一个未访问格点。因此归纳成立。旋转和反射保持“第一个未访问格点”的关系，覆盖其余 $21$ 种排列。注意 \tt{NESW} 的四个阶段沿不同方向伸展，不能把整块轨迹当作同一个平移；但\textbf{固定阶段}的坐标总是等差数列。

对任意排列，设 $v_i$ 是执行 $i$ 次指令后的坐标，$n=4q+r$，$0\le r<4$。当 $q\ge2$ 时，$v_{4q+r}=v_{8+r}+(q-2)(v_{12+r}-v_{8+r})$。较小的 $n$ 直接读取模拟结果。坐标与乘积需用 64 位整数。

每组测试只模拟常数次指令，时间 $O(T)$，额外空间 $O(1)$。
